\partsection{B.4}{Verification: commands, results, AC to test, excerpts, TDD}\label{B:4}

\subsubsection*{Commands, tested commit and counts}

Every local command below ran on 2026-11-25 on the tree of \texttt{v0.1.1.0} (\texttt{3a91ea4}, a clean worktree at that commit) on the author's Windows machine with Docker Desktop (PostgreSQL 17 containers through \path{internal/pgtest}). CI ran the same suites on \texttt{3a91ea4} itself: CI run 45614889848 (push to \texttt{main}), every job \texttt{success} (\texttt{image}, \texttt{code\_quality}, \texttt{test}, \texttt{build}, \texttt{api}, \texttt{release}), followed by Deploy demo run 45615736377 (Part E).
\todo{recapture on candidate: rerun the commands on the candidate commit and replace the counts and run IDs}

\textbf{Go API: tests.} \texttt{bun run go-test 02} runs \texttt{go test -json -count=1 ./...} in \path{apps/api} and keeps the JSON stream that feeds the AC table.

\begin{shell}
$ cd apps/api && go test -json -count=1 ./...        (via: bun run go-test 02)
packages with tests: 5   (apps/api, internal/auth, internal/claim, internal/money, migrations)
top-level tests: 178 passed, 0 failed, 0 skipped
subtests:        241 passed, 0 failed, 0 skipped
per package (top-level tests passed): apps/api 114, internal/claim 47, internal/auth 11,
                                      internal/money 3, migrations 3
exit status 0   (6 packages have no test files: env, internal/api, internal/db,
                 internal/httperr, internal/pgtest, internal/system)
\end{shell}

\textbf{Go API: race detector and the 100\% coverage gate.}

\begin{shell}
$ cd apps/api && bun run test:coverage
  = go test -race -coverpkg=./... -coverprofile=coverage.out ./... && go-test-coverage --config=.testcoverage.yml
ok  apps/api                       196.141s
ok  apps/api/internal/auth           2.267s
ok  apps/api/internal/claim        110.250s
ok  apps/api/internal/money          5.935s
ok  apps/api/migrations             35.805s
File coverage threshold (100%) satisfied:     PASS
Package coverage threshold (100%) satisfied:  PASS
Total coverage threshold (100%) satisfied:    PASS
Total test coverage: 100% (956/956)           exit status 0, no data race reported
\end{shell}

\textbf{Web (Vitest), from CI.} The full web suite and its coverage are not run on the author's machine: a full local run times out under load, so CI is the verdict (CI run 45614889848, job \texttt{test}, \texttt{bun run test:coverage}).

\begin{shell}
CI run 45614889848 on 3a91ea4, job test:
 Test Files  22 passed (22)
      Tests  182 passed (182)
 Statements : 100% ( 537/537 )    Branches : 100% ( 313/313 )
 Functions  : 100% ( 182/182 )    Lines    : 100% ( 528/528 )
CI job api: golangci-lint 0 issues; steps Check generated code, Check modules are tidy and
            Check vulnerabilities passed;
            Total test coverage: 100% (956/956)
\end{shell}

\textbf{End-to-end (Playwright, local seeded stack, Chrome).} The specs run on a local stack before each release, not in CI (Product DoD 0.3 item 4). Playwright's own Chromium was not installed on this machine, so the run used the installed Google Chrome through a temporary Playwright config that was not committed; the specs are unchanged.

\begin{shell}
$ SEED_USER_PASSWORD=<set in the process environment> playwright test --timeout=60000   (apps/web/e2e)
  F01-TEST-24, F02-TEST-21, F03-TEST-12, F04-TEST-23, F04-TEST-24,
  F05-TEST-25, F06-TEST-15, F07-TEST-12
  8 passed (1.4m)
\end{shell}

\textbf{Failed and skipped.} Go: 0 failed, 0 skipped (178 top-level tests, 241 subtests). Web in CI: 182 of 182. E2E: 8 of 8. No test was disabled, skipped or weakened to obtain these results. The one local-only web failure recorded in the contracts (\texttt{F01-TEST-20} in \path{auth-routing.test.tsx}, failing alone on the author's machine and passing in CI) is listed in the back matter.

% Generated by scripts/gen-ac-table.ts at 3a91ea4. Do not edit.

\subsubsection*{Automated tests (56 planned tests with an actual test)}
{\footnotesize\setlength{\tabcolsep}{3pt}
\begin{longtable}{@{}p{1.9cm}p{2.3cm}p{1.8cm}p{6.6cm}p{3.6cm}@{}}
\toprule \textbf{Planned test} & \textbf{AC} & \textbf{Level} & \textbf{Actual test (file, name)} & \textbf{Result} \\ \midrule \endhead
F05-TEST-01 & F05-AC-01 & Integration & \path{api/claims_drafts_test.go}\newline \hspace*{1em}\seqsplit{TestF05TEST01SaveDraftWithoutSubmission} & pass (1 run) \\
F05-TEST-02 & F05-AC-02 & Integration (module) & \path{api/internal/claim/drafts_test.go}\newline \hspace*{1em}\seqsplit{TestF05TEST02DraftsCheckItemAndTotalLimitsButSkipDates} & pass (1 run) \\
F05-TEST-03 & F05-AC-03 & Integration & \path{api/claims_drafts_test.go}\newline \hspace*{1em}\seqsplit{TestF05TEST03FinanceOfficerCannotSaveDraft} & pass (1 run) \\
F05-TEST-04 & F05-AC-03 & Integration (module) & \path{api/internal/claim/drafts_test.go}\newline \hspace*{1em}\seqsplit{TestF05TEST04CreateDraftRequiresCurrentReviewer} & pass (1 run) \\
F05-TEST-05 & F05-AC-04 & Integration & \path{api/claims_drafts_test.go}\newline \hspace*{1em}\seqsplit{TestF05TEST05OnlyOwnerCanOpenDraft} & pass (1 run) \\
F05-TEST-06 & F05-AC-04, F05-AC-06, F05-AC-11 & Unit & \path{api/internal/claim/policy_test.go}\newline \hspace*{1em}\seqsplit{TestF05TEST06DraftAndReturnedPolicy} & pass (1 run) \\
F05-TEST-07 & F05-AC-05 & Integration & \path{api/claims_drafts_test.go}\newline \hspace*{1em}\seqsplit{TestF05TEST07EditReplacesItemsAndRecalculatesTotal} & pass (1 run) \\
F05-TEST-08 & F05-AC-05 & Integration (module) & \path{api/internal/claim/drafts_test.go}\newline \hspace*{1em}\seqsplit{TestF05TEST08InvalidEditLeavesDraftUnchanged} & pass (1 run) \\
F05-TEST-09 & F05-AC-06 & Integration & \path{api/claims_drafts_test.go}\newline \hspace*{1em}\seqsplit{TestF05TEST09EditRequiresOwnerAndEditableStatus} & pass (1 run) \\
F05-TEST-10 & F05-AC-07 & Integration & \path{api/claims_drafts_test.go}\newline \hspace*{1em}\seqsplit{TestF05TEST10DraftMutationsRequireCurrentIfMatch} & pass (1 run) \\
F05-TEST-11 & F05-AC-08 & Integration & \path{api/claims_drafts_test.go}\newline \hspace*{1em}\seqsplit{TestF05TEST11SubmitDraftRecordsFirstSubmissionAndQueuesClaim} & pass (1 run) \\
F05-TEST-12 & F05-AC-09 & Integration (module) & \path{api/internal/claim/drafts_test.go}\newline \hspace*{1em}\seqsplit{TestF05TEST12SubmissionChecksExpenseDateWindow} & pass (1 run) \\
F05-TEST-13 & F05-AC-09 & Integration & \path{api/claims_drafts_test.go}\newline \hspace*{1em}\seqsplit{TestF05TEST13SubmitDraftReportsExpenseDateField} & pass (1 run) \\
F05-TEST-14 & F05-AC-10 & Integration (module) & \path{api/internal/claim/drafts_test.go}\newline \hspace*{1em}\seqsplit{TestF05TEST14DraftSubmissionUsesOwnersCurrentReviewer} & pass (1 run) \\
F05-TEST-15 & F05-AC-11 & Integration (module) & \path{api/internal/claim/drafts_test.go}\newline \hspace*{1em}\seqsplit{TestF05TEST15DraftSubmissionRequiresOwnersCurrentReviewer} & pass (1 run) \\
F05-TEST-16 & F05-AC-12 & Integration & \path{api/claims_drafts_test.go}\newline \hspace*{1em}\seqsplit{TestF05TEST16OnlyDraftOrReturnedCanBeSubmitted} & pass (1 run) \\
F05-TEST-17 & F05-AC-12 & Integration (module) & \path{api/internal/claim/drafts_test.go}\newline \hspace*{1em}\seqsplit{TestF05TEST17ConcurrentDraftSubmissionsOneWins} & pass (1 run) \\
F05-TEST-18 & F05-AC-13 & Integration (module) & \path{api/internal/claim/drafts_test.go}\newline \hspace*{1em}\seqsplit{TestF05TEST18HistoryFailureRollsBackDraftSubmission} & pass (1 run) \\
F05-TEST-19 & F05-AC-14 & Integration & \path{api/claims_drafts_test.go}\newline \hspace*{1em}\seqsplit{TestF05TEST19DeleteRequiresOwnerAndDraftStatus} & pass (1 run) \\
F05-TEST-20 & F05-AC-15 & Integration & \path{api/claims_drafts_test.go}\newline \hspace*{1em}\seqsplit{TestF05TEST20HistoryPreventsDirectClaimDeletion} & pass (1 run) \\
F05-TEST-21 & F05-AC-16 & Component & \path{web/src/features/claims/__test__/draft-workflow.test.tsx}\newline \hspace*{1em}F05-TEST-21: keeps input on a draft save error, disables both actions while saving, and opens the draft & pass (Vitest, CI at 3a91ea4) \\
F05-TEST-22 & F05-AC-17 & Component & \path{web/src/features/claims/__test__/draft-workflow.test.tsx}\newline \hspace*{1em}F05-TEST-22: shows draft placeholders and hides actions absent from policy\newline \hspace*{1em}F05-TEST-22: shows each invalid line, then submits with the version and refetches\newline \hspace*{1em}F05-TEST-22: reloads on conflict and confirms before deleting\newline \hspace*{1em}F05-TEST-22: shows a claim-wide validation message without an unavailable edit action\newline \hspace*{1em}F05-TEST-22: reloads after a rejected transition and shows \$message\newline \hspace*{1em}F05-TEST-22: reports a failed submission when \$name\newline \hspace*{1em}F05-TEST-22: closes the delete confirmation when \$name & pass (Vitest, CI at 3a91ea4) \\
F05-TEST-23 & F05-AC-18 & Component & \path{web/src/features/claims/__test__/draft-workflow.test.tsx}\newline \hspace*{1em}F05-TEST-23: prefills the form, preserves invalid input, and reloads on conflict\newline \hspace*{1em}F05-TEST-23: declines the edit route without the edit action\newline \hspace*{1em}F05-TEST-23: shows the access result when saving returns \$status\newline \hspace*{1em}F05-TEST-23: returns to detail after a rejected edit and shows \$message\newline \hspace*{1em}F05-TEST-23: shows no form error while a conflicting draft reloads\newline \hspace*{1em}F05-TEST-23: retries a failed conflict reload and opens the latest draft & pass (Vitest, CI at 3a91ea4) \\
F05-TEST-24 & F05-AC-19 & Component & \path{web/src/features/claims/__test__/draft-workflow.test.tsx}\newline \hspace*{1em}F05-TEST-24: lists a draft with a neutral badge and a dash in Submitted & pass (Vitest, CI at 3a91ea4) \\
F05-TEST-25 & F05-AC-20 & E2E & \path{web/e2e/f05-drafts.spec.ts}\newline \hspace*{1em}F05-TEST-25: save, edit, submit, review and delete drafts & pass (E2E on a local stack at 3a91ea4, not in CI) \\
F05-TEST-26 & F04-AC-19 (seed) & Integration & \path{api/seed_test.go}\newline \hspace*{1em}\seqsplit{TestF04TEST22SeedAddsSampleClaimsOnce} & pass (1 run) \\
F06-TEST-01 & F06-AC-01 & Integration & \path{api/claims_return_test.go}\newline \hspace*{1em}\seqsplit{TestF06TEST01ReviewerReturnsASubmittedClaim} & pass (1 run) \\
F06-TEST-02 & F06-AC-02 & Integration (module) & \path{api/internal/claim/return_test.go}\newline \hspace*{1em}\seqsplit{TestF06TEST02ReturnReasonFollowsTheRejectRuleBR04} & pass (1 run) \\
F06-TEST-03 & F06-AC-02 & Integration & \path{api/claims_return_test.go}\newline \hspace*{1em}\seqsplit{TestF06TEST03ReturnWithoutAReasonIs422} & pass (1 run) \\
F06-TEST-04 & F06-AC-03 & Integration & \path{api/claims_return_test.go}\newline \hspace*{1em}\seqsplit{TestF06TEST04OnlyTheAssignedReviewerCanReturn} & pass (1 run) \\
F06-TEST-05 & F06-AC-04 & Integration & \path{api/claims_return_test.go}\newline \hspace*{1em}\seqsplit{TestF06TEST05OnlyASubmittedClaimCanBeReturned} & pass (1 run) \\
F06-TEST-06 & F06-AC-05 & Unit & \path{api/internal/claim/policy_test.go}\newline \hspace*{1em}\seqsplit{TestF06TEST06ReturnIsAllowedBelowTheReturnLimit} & pass (1 run) \\
F06-TEST-07 & F06-AC-05 & Integration (module) & \path{api/claims_return_test.go}\newline \hspace*{1em}\seqsplit{TestF06TEST07ReturnLimitIs409AndApproveStillWorks}\newline \path{api/internal/claim/return_test.go}\newline \hspace*{1em}\seqsplit{TestF06TEST07FourthReturnIsRefusedAndApproveStillWorks} & pass (2 run) \\
F06-TEST-08 & F06-AC-06 & Integration & \path{api/claims_return_test.go}\newline \hspace*{1em}\seqsplit{TestF06TEST08ReturnedClaimLeavesTheQueueUntilResubmitted} & pass (1 run) \\
F06-TEST-09 & F06-AC-07 & Integration & \path{api/claims_resubmit_test.go}\newline \hspace*{1em}\seqsplit{TestF06TEST09OwnerResubmitsToTheSameReviewer} & pass (1 run) \\
F06-TEST-10 & F06-AC-08 & Integration (module) & \path{api/internal/claim/resubmit_test.go}\newline \hspace*{1em}\seqsplit{TestF06TEST10ResubmissionCountsNinetyDaysFromTheFirstSubmission} & pass (1 run) \\
F06-TEST-11 & F06-AC-09 & Integration & \path{api/claims_resubmit_test.go}\newline \hspace*{1em}\seqsplit{TestF06TEST11OnlyTheOwnerEditsAndResubmitsAReturnedClaim} & pass (1 run) \\
F06-TEST-12 & F06-AC-10 & Integration (module) & \path{api/internal/claim/resubmit_test.go}\newline \hspace*{1em}\seqsplit{TestF06TEST12HistoryFailureRollsBackTheResubmission}\newline \path{api/internal/claim/return_test.go}\newline \hspace*{1em}\seqsplit{TestF06TEST12HistoryFailureRollsBackTheReturn} & pass (2 run) \\
F06-TEST-13 & F06-AC-11 & Component & \path{web/src/features/reviews/__test__/return-panel.test.tsx}\newline \hspace*{1em}F06-TEST-13: shows Return with the number of returns used\newline \hspace*{1em}F06-TEST-13: shows no Return button when return is not allowed\newline \hspace*{1em}F06-TEST-13: with no returns left it says to approve or reject\newline \hspace*{1em}F06-TEST-13: the dialog counts the reason, and a 422 puts the message under the field with the text kept\newline \hspace*{1em}F06-TEST-13: return requires 10 characters, as reject does\newline \hspace*{1em}F06-TEST-13: the buttons are disabled with a spinner while the return is sent\newline \hspace*{1em}F06-TEST-13: a return shows the toast and goes back to the review queue\newline \hspace*{1em}F06-TEST-13: a 409 shows the changed toast, closes the dialog, and reads the claim again & pass (Vitest, CI at 3a91ea4) \\
F06-TEST-14 & F06-AC-12 & Component & \path{web/src/features/claims/__test__/resubmit.test.tsx}\newline \hspace*{1em}F06-TEST-14 (Edit, Resubmit, toast): the owner of a returned claim sees Edit and Resubmit, not Submit claim, and resubmitting shows the toast and refetches\newline \hspace*{1em}F06-TEST-14 (422): lists each broken line as Line n and keeps the claim returned\newline \path{web/src/features/claims/__test__/return-banner.test.tsx}\newline \hspace*{1em}F06-TEST-14 (banner and count; Edit and Resubmit are \#44): shows the latest return reason with who returned it and when, and the number of returns\newline \hspace*{1em}F06-TEST-14: shows the banner without a reason when the return entry has none & pass (Vitest, CI at 3a91ea4) \\
F06-TEST-15 & F06-AC-13 & E2E & \path{web/e2e/f06-return-resubmit.spec.ts}\newline \hspace*{1em}F06-TEST-15: req.anan submits, rev.malee returns with a reason, req.anan edits and resubmits, rev.malee approves, and the history reads Submit, Return, Resubmit, Approve & pass (E2E on a local stack at 3a91ea4, not in CI) \\
F06-TEST-16 & F04-AC-19 (seed) & Integration & \path{api/seed_test.go}\newline \hspace*{1em}\seqsplit{TestF04TEST22SeedAddsSampleClaimsOnce} & pass (1 run) \\
F06-TEST-17 & F06-AC-07 & Integration (module) & \path{api/internal/claim/resubmit_test.go}\newline \hspace*{1em}\seqsplit{TestF06TEST17ResubmissionWorksWithoutAnOwnerReviewer} & pass (1 run) \\
F07-TEST-01 & F07-AC-01 & Integration & \path{api/claims_withdraw_test.go}\newline \hspace*{1em}\seqsplit{TestF07TEST01OwnerWithdrawsASubmittedOrReturnedClaim} & pass (1 run) \\
F07-TEST-02 & F07-AC-02 & Integration & \path{api/claims_withdraw_test.go}\newline \hspace*{1em}\seqsplit{TestF07TEST02WithdrawCommentIsOptionalWithAtMost500Characters} & pass (1 run) \\
F07-TEST-03 & F07-AC-03 & Integration & \path{api/claims_withdraw_test.go}\newline \hspace*{1em}\seqsplit{TestF07TEST03OnlyASubmittedOrReturnedClaimCanBeWithdrawnBR16} & pass (1 run) \\
F07-TEST-04 & F07-AC-04 & Integration & \path{api/claims_withdraw_test.go}\newline \hspace*{1em}\seqsplit{TestF07TEST04OnlyTheOwnerCanWithdraw} & pass (1 run) \\
F07-TEST-05 & F07-AC-04, F07-AC-06 & Unit & \path{api/internal/claim/policy_test.go}\newline \hspace*{1em}\seqsplit{TestF07TEST05WithdrawIsTheOwnersOnSubmittedOrReturned} & pass (1 run) \\
F07-TEST-06 & F07-AC-05 & Integration & \path{api/claims_withdraw_test.go}\newline \hspace*{1em}\seqsplit{TestF07TEST06WithdrawNeedsTheCurrentIfMatch} & pass (1 run) \\
F07-TEST-07 & F07-AC-06 & Integration & \path{api/claims_withdraw_test.go}\newline \hspace*{1em}\seqsplit{TestF07TEST07AWithdrawnClaimCannotBeActedOnAgain} & pass (1 run) \\
F07-TEST-08 & F07-AC-06 & Integration (module) & \path{api/internal/claim/withdraw_test.go}\newline \hspace*{1em}\seqsplit{TestF07TEST08ConcurrentWithdrawAndApproveHaveOneWinner} & pass (1 run) \\
F07-TEST-09 & F07-AC-07 & Integration (module) & \path{api/internal/claim/withdraw_test.go}\newline \hspace*{1em}\seqsplit{TestF07TEST09HistoryFailureRollsBackTheWithdrawal} & pass (1 run) \\
F07-TEST-10 & F07-AC-08 & Integration & \path{api/claims_withdraw_test.go}\newline \hspace*{1em}\seqsplit{TestF07TEST10AWithdrawnClaimLeavesTheReviewQueueButStaysReadable} & pass (1 run) \\
F07-TEST-11 & F07-AC-09 & Component & \path{web/src/features/claims/__test__/withdraw-claim.test.tsx}\newline \hspace*{1em}F07-TEST-11: shows Withdraw only when the policy gives it\newline \hspace*{1em}F07-TEST-11: opens a dialog with the optional comment, a 0/500 counter and a limit of 500 characters\newline \hspace*{1em}F07-TEST-11: disables the confirm button while sending, then shows the toast, the Withdrawn badge and no actions\newline \hspace*{1em}F07-TEST-11: after the withdrawal it shows the toast, the Withdrawn badge and an empty action bar\newline \hspace*{1em}F07-TEST-11: a 409 CONFLICT shows the reload toast and reads the claim again\newline \hspace*{1em}F07-TEST-11: a 409 INVALID\_TRANSITION shows its message and reads the claim again\newline \hspace*{1em}F07-TEST-11: Cancel and Escape close the dialog without sending anything\newline \hspace*{1em}F07-TEST-11: Escape does not close the dialog while the withdrawal is sending\newline \hspace*{1em}F07-TEST-11: a failure that is not a problem document shows the generic toast and closes the dialog & pass (Vitest, CI at 3a91ea4) \\
F07-TEST-12 & F07-AC-10 & E2E & \path{web/e2e/f07-withdraw.spec.ts}\newline \hspace*{1em}F07-TEST-12: req.anan submits and withdraws a claim with a comment, it shows Withdrawn in My Claims, and rev.malee & pass (E2E on a local stack at 3a91ea4, not in CI) \\
F07-TEST-13 & F04-AC-19 (seed) & Integration & \path{api/seed_test.go}\newline \hspace*{1em}\seqsplit{TestF04TEST22SeedAddsSampleClaimsOnce} & pass (1 run) \\
\bottomrule\end{longtable}}

\subsubsection*{Manual and CI-recorded checks (0)}
{\footnotesize\setlength{\tabcolsep}{3pt}
\begin{longtable}{@{}p{1.9cm}p{2.3cm}p{1.8cm}p{10.2cm}@{}}
\toprule \textbf{Planned test} & \textbf{AC} & \textbf{Level} & \textbf{Status} \\ \midrule \endhead
\bottomrule\end{longtable}}

\subsubsection*{Planned but not implemented yet (0)}
{\footnotesize\setlength{\tabcolsep}{3pt}
\begin{longtable}{@{}p{1.9cm}p{4.2cm}p{2.2cm}p{7.9cm}@{}}
\toprule \textbf{Planned test} & \textbf{AC} & \textbf{Level} & \textbf{Result} \\ \midrule \endhead
\bottomrule\end{longtable}}


\noindent The Result column shows the result of each automated test on \texttt{3a91ea4}: Go tests from the \texttt{go test -json} run above, web component tests from CI run 45614889848, E2E tests from the local Playwright run. The seed rows (F05-TEST-26, F06-TEST-16, F07-TEST-13) are found through \path{docs/traceability.md} and are subtests of \texttt{TestF04TEST22SeedAddsSampleClaimsOnce}; the pass shown is that of the top-level test. There are no manual checks and no planned test without an actual test in Sprint 2.

\subsubsection*{Test excerpts}

\textbf{Domain rule: BR-14, a claim is returned at most three times} (F06-TEST-07, F06-AC-05).
\path{apps/api/internal/claim/return_test.go}, commit \texttt{3a91ea4}, lines 98--132.
The setup is the seeded requester and reviewer and the real \texttt{claim.Service} on PostgreSQL; the helper \texttt{returnAndResubmit} returns and resubmits the claim \texttt{MaxReturns} times; the assertions are that the fourth return fails with \texttt{ErrInvalidTransition} for the return limit, that the stored status, return count and version are unchanged, and that Approve still works.

\VerbatimInput[tabsize=4, fontsize=\scriptsize, frame=single, framesep=5pt, rulecolor=\color{black!25}, breaklines=true]{content/data/excerpt-br14.go}

\textbf{Access restriction: only the owner can withdraw} (F07-TEST-04, F07-AC-04, NFR-02).
\path{apps/api/claims_withdraw_test.go}, commit \texttt{3a91ea4}, lines 147--169.
The setup is the real router with the seeded users: \texttt{req.anan} submits a claim; the assigned reviewer \texttt{rev.malee} and another requester \texttt{req.ploy} each send \texttt{POST /claims/\{id\}/withdraw} with the current version. The assertions are the 403 \texttt{FORBIDDEN} problem, that the body leaks none of \texttt{data}, \texttt{claimNo} or the title, and that the stored claim is unchanged.

\VerbatimInput[tabsize=4, fontsize=\scriptsize, frame=single, framesep=5pt, rulecolor=\color{black!25}, breaklines=true]{content/data/excerpt-f07-access.go}

\textbf{Execution results of the two excerpts} (run on \texttt{3a91ea4}):

\VerbatimInput[tabsize=4, fontsize=\scriptsize, frame=single, framesep=5pt, rulecolor=\color{black!25}, breaklines=true]{content/data/excerpt-results.txt}

\subsubsection*{TDD: the failing test first, then the code}

Two pairs, each a test commit that fails and the commit that makes it pass. All four are commits of PR \#108 (50 commits on \path{feat/f06-f07-return-resubmit-withdraw}); the PR was squash-merged into \texttt{sprint-2} as \texttt{aaf17af}, so they stay reachable under the PR's Commits tab and as \texttt{refs/pull/108/head}. Author of all four: \texttt{Kiatisakk}. Output: the same command run on a clean checkout of each commit (Docker Desktop running); container start-up lines removed.

\begin{center}\small
\begin{tabular}{@{}lllp{7.7cm}@{}}
\toprule \textbf{Rule, test} & \textbf{Step, commit} & \textbf{Author date (UTC+7)} & \textbf{Subject} \\ \midrule
BR-14, F06-TEST-07 & Red \texttt{a547ac0} & 2026-11-25 01:13:33 & \path{test: refuse a fourth return of a claim (F06-TEST-07)} \\
 & Green \texttt{71a8e43} & 2026-11-25 01:17:19 & \path{feat: limit a claim to three returns} \\
BR-16, F07-TEST-03 & Red \texttt{53e8827} & 2026-11-25 02:24:32 & \path{test: only a submitted or returned claim can be withdrawn (F07-TEST-03)} \\
 & Green \texttt{38e0165} & 2026-11-25 02:35:27 & \path{feat: add the withdraw operation for a submitted or returned claim (F07-TEST-03)} \\
\bottomrule
\end{tabular}
\end{center}

\textbf{BR-14: why the red test failed.} At \texttt{a547ac0} the commit adds the test and only the names it needs to compile (the constant \texttt{MaxReturns} and the error \texttt{ErrReturnLimit}); the service does not apply the limit. The fourth return therefore reached PostgreSQL, whose \texttt{CHECK (return\_count <= 3)} refused it, and the service returned a database error where the test expects \texttt{ErrInvalidTransition}. At \texttt{71a8e43} the service checks the count before the update and the test passes.

\textbf{Red, \texttt{a547ac0}:}

\VerbatimInput[tabsize=4, fontsize=\scriptsize, frame=single, framesep=5pt, rulecolor=\color{black!25}, breaklines=true]{content/data/tdd-br14-red.txt}

\textbf{Green, \texttt{71a8e43}:}

\VerbatimInput[tabsize=4, fontsize=\scriptsize, frame=single, framesep=5pt, rulecolor=\color{black!25}, breaklines=true]{content/data/tdd-br14-green.txt}

\textbf{BR-16: why the red test failed.} At \texttt{53e8827} only the test exists: the withdraw route is not in the API, so each of the four starting statuses (\texttt{DRAFT}, \texttt{APPROVED}, \texttt{REJECTED}, \texttt{WITHDRAWN}) got 404 \texttt{NOT\_FOUND} where the test expects 409. At \texttt{38e0165} the operation and its lifecycle rule exist, and all four cases return 409 \texttt{INVALID\_TRANSITION} with BR-16.

\textbf{Red, \texttt{53e8827}:}

\VerbatimInput[tabsize=4, fontsize=\scriptsize, frame=single, framesep=5pt, rulecolor=\color{black!25}, breaklines=true]{content/data/tdd-br16-red.txt}

\textbf{Green, \texttt{38e0165}:}

\VerbatimInput[tabsize=4, fontsize=\scriptsize, frame=single, framesep=5pt, rulecolor=\color{black!25}, breaklines=true]{content/data/tdd-br16-green.txt}

\textbf{Other test-first pairs and honest exceptions.} The Test column of the F-05, F-06 and F-07 result tables in \path{docs/traceability.md} (Part D) lists a red and a green commit for every other test that went red first, for example BR-10 (F05-TEST-09, 19: red \texttt{ec2179d}, green \texttt{dd31979}) and BR-07 (F05-TEST-12: red \texttt{4dcf8d4}, green \texttt{06c5fa0}) in PR \#104. Where a test could not go red first the same table says so and the contract accepts it: F06-TEST-03, 04, 06, 08, 09, 10, 11, 12, 17 and F07-TEST-04, 06 to 10 passed when written because earlier green commits already covered them; F06-TEST-15 and F07-TEST-12 were written after the code; the BR-07 resubmission window (F06-TEST-10) has a mutation check instead of a red commit, accepted in F-06 contract revision 0.3 after review on PR \#108.
